\begin{abstract}
Spectral analyzers in audio hosts run on the engine thread even when they
produce only display data, so a frame's windowing, transform, and magnitude
processing can land in a single sample call. We present an exact-hop execution
contract that spreads the complete analysis---window preparation, FFT
butterflies, magnitude processing, smoothing, and output---over the sample
calls of one hop. The contract bounds the logical work in each call, retains
the selected input while capture continues, publishes on a frame clock that is
independent of arithmetic completion, and transfers complete spectra through
an owned mailbox. A descriptive study of 776 fresh processes in two sessions
on one Apple M1 Pro separates work placement from native FFT-kernel choice and
measures complete modules inside VCV Rack's engine. For scalar
$N=4096$, $H=1024$ analysis in 64-sample blocks, distributing a fixed
implementation reduces the median per-hop peak block time from
\StudyBatchPeak{} to \StudyCorePeak\,$\mu$s at \StudyPlacementCostRatio{}
times the instrumented aggregate cost. With four aligned four-channel modules
on one engine thread, the production pipeline, a native vDSP batch path, and a
native hybrid have typical hop peaks of \StudyHostCorePeak,
\StudyHostBatchPeak, and \StudyHostHybridPeak\,$\mu$s---%
\StudyHostCoreBudgetPercent\%, \StudyHostBatchBudgetPercent\%, and
\StudyHostHybridBudgetPercent\% of the block budget---while the two
distributed paths add about 30\,ms of logical publication age. Native hybrids
and shorter completion horizons offer intermediate tradeoffs. Multithreaded
engine overhead and release lateness, however, keep lower burst peaks from
translating reliably into fewer missed deadlines. The result is an
analysis-specific scheduling contract and a measured basis for choosing
execution granularity, not a universal performance ranking.
\end{abstract}
